% =================================================
% 文件: 05_假负载.tex
% 章节: 第21章 假负载
% 所属部分: 仪器仪表
% 版本: v1.0
% =================================================

\chapter{假负载}
\label{chap:dummy_load}

假负载（Dummy Load）是一种用于测试发射机的专用设备，用于替代天线作为负载。

从电路角度看，天线对发射机而言只是一个具有特定阻抗的负载：在设计频率附近，调谐良好的天馈系统在馈线输入端呈现接近纯电阻的阻抗，其值通常为 $50\,\Omega$（部分广播与电视系统为 $75\,\Omega$）。发射机末级正是按这一阻抗设计其匹配网络与工作状态的。假负载的作用就是提供一个与天馈系统等效、但不向空间辐射能量的负载，使发射机能够在完全正常的负载条件下工作，而测试者可以安全、可重复地测量其各项指标。

使用假负载的必要性首先来自器件安全。发射机末级若在不接负载（开路）或负载短路的情况下加激励，功率无法送出，全部或大部分能量将在末级器件与匹配网络内部消耗，并在开路点形成很高的射频电压。对电子管发射机，失配会造成屏极耗散剧增、屏极发红甚至打火；对晶体管功率放大器，失配会使集电极（或漏极）电压峰值远超正常值，极易造成二次击穿而瞬间损坏。即便器件未立即损坏，长期在高驻波下工作也会加速老化、引起自激振荡和寄生辐射。其次是测试的规范性：只有在标准负载下测得的功率、频率、调制与杂散指标才具有可比性；用实际天线测量时，天线阻抗随环境、高度和邻近物体变化，数据无法复现。最后是电磁环境要求：室内调试、维修间检修以及生产线批量测试都不允许持续向空间辐射大功率信号，假负载把能量就地转换为热能，从根本上避免了对其他设备和无线电业务的干扰。

\section{工作原理}
\label{sec:dummy_load_principle}

假负载通过电阻元件吸收发射机输出的全部功率，并将其转换为热能散发，从而模拟天线的负载特性，同时避免信号辐射到空间。

理想的假负载在整个工作频段内应等效为一个数值恒定的纯电阻，即其输入阻抗
\[
Z_{in} = R + \mathrm{j}X, \qquad R = Z_0,\quad X = 0 .
\]
实际器件不可能完全满足这一条件：电阻体存在分布电容，引线和结构件存在寄生电感，随着频率升高，电抗分量 $X$ 逐渐增大，阻抗偏离 $Z_0$，性能随之变差。因此频率范围与驻波比总是成对给出的指标，脱离频率谈驻波比没有意义。

负载与馈线不匹配的程度用反射系数描述
\[
\Gamma = \frac{Z_L - Z_0}{Z_L + Z_0},
\]
式中 $Z_L$ 为负载阻抗，$Z_0$ 为传输线特性阻抗。$\Gamma$ 一般是复数，其模 $|\Gamma|$ 表示反射波与入射波的幅度比。工程上更常用电压驻波比表示同一事实
\[
VSWR = \frac{U_{\max}}{U_{\min}} = \frac{1 + |\Gamma|}{1 - |\Gamma|},
\qquad |\Gamma| = \frac{VSWR - 1}{VSWR + 1},
\]
以及以分贝表示的回波损耗
\[
RL = -20\lg|\Gamma| \quad (\mathrm{dB}).
\]
入射功率中被反射回去的比例为 $|\Gamma|^2$，被负载吸收的比例为 $1 - |\Gamma|^2$，由此定义失配损耗
\[
L_{mis} = -10\lg\left(1 - |\Gamma|^{2}\right) \quad (\mathrm{dB}).
\]
下表给出几个常用驻波比对应的反射系数、回波损耗与反射功率百分比，便于估算。可以看出，$VSWR = 1.5$ 时反射功率只有 $4.0\%$，对发射机而言影响有限；但当 $VSWR$ 达到 $3.0$ 时，已有 $25\%$ 的功率返回末级，多数固态发射机的保护电路会在此之前就动作降功率或关机。

\begin{table}[htbp]
\centering
\caption{驻波比、反射系数、回波损耗与反射功率的对应关系}
\begin{tabular}{|p{2.2cm}|p{2.6cm}|p{3.0cm}|p{3.4cm}|}
\hline
\textbf{VSWR} & \textbf{$|\Gamma|$} & \textbf{回波损耗 (dB)} & \textbf{反射功率百分比} \\
\hline
1.00 & 0.000 & $\infty$ & $0\%$ \\
\hline
1.05 & 0.024 & 32.3 & $0.06\%$ \\
\hline
1.10 & 0.048 & 26.4 & $0.23\%$ \\
\hline
1.20 & 0.091 & 20.8 & $0.83\%$ \\
\hline
1.50 & 0.200 & 14.0 & $4.0\%$ \\
\hline
2.00 & 0.333 & 9.5 & $11.1\%$ \\
\hline
3.00 & 0.500 & 6.0 & $25.0\%$ \\
\hline
\end{tabular}
\end{table}

假负载吸收的功率全部变为热量，其电与热的关系是选型和自制时最需要定量把握的部分。在阻抗匹配的条件下，负载上的射频电压有效值 $U$ 与吸收功率 $P$ 满足
\[
P = \frac{U^2}{R}, \qquad U = \sqrt{P R},
\]
对 $R = 50\,\Omega$ 的系统，还有峰值电压 $U_p = \sqrt{2}\,U = \sqrt{2PR}$。功率也常以 dBm 表示，即以 1~mW 为参考的分贝值
\[
P(\mathrm{dBm}) = 10\lg\frac{P}{1\ \mathrm{mW}} .
\]
下表列出 $50\,\Omega$ 系统中几个典型功率对应的电压与 dBm 值。这些数值提醒使用者：即使是百瓦级的普通电台，负载端的射频电压峰值已达百伏量级，千瓦级设备更超过 300~V，因此假负载的绝缘、接头耐压与操作安全都不容忽视，绝不可在发射状态下触碰负载或拆卸接头。

\begin{table}[htbp]
\centering
\caption{$50\,\Omega$ 系统中功率与电压、dBm 的换算}
\begin{tabular}{|p{2.4cm}|p{2.8cm}|p{2.8cm}|p{2.6cm}|}
\hline
\textbf{功率 $P$} & \textbf{有效值 $U$ (V)} & \textbf{峰值 $U_p$ (V)} & \textbf{电平 (dBm)} \\
\hline
1 mW & 0.224 & 0.316 & 0 \\
\hline
1 W & 7.07 & 10.0 & 30 \\
\hline
10 W & 22.4 & 31.6 & 40 \\
\hline
100 W & 70.7 & 100 & 50 \\
\hline
1000 W & 224 & 316 & 60 \\
\hline
\end{tabular}
\end{table}

热的一侧则由热阻决定。设负载吸收功率为 $P$，从电阻体（结）到环境的总热阻为 $R_{th}$（单位 $^\circ\mathrm{C/W}$），环境温度为 $T_a$，则稳态时电阻体温度近似为
\[
T_j = T_a + P \cdot R_{th},
\]
总热阻为各段串联之和，即电阻体到外壳、外壳到散热器、散热器到空气三段之和
\[
R_{th} = R_{th(j\text{-}c)} + R_{th(c\text{-}s)} + R_{th(s\text{-}a)} .
\]
这一简单关系说明了假负载散热设计的全部要点：欲在给定功率下降低温升，只能减小热阻——增大散热面积、改善接触（涂导热硅脂、加大压紧力）、加装风扇强迫风冷，或改用油浸、水冷。同时也说明了额定功率必须与环境温度和工作时间联系起来看：手册标称的连续功率通常是在规定环境温度和规定散热条件（例如带风扇）下的值，环境温度升高或去掉风扇后必须降额使用。

\section{主要特性}
\label{sec:dummy_load_features}
\begin{itemize}
    \item \textbf{阻抗匹配}：通常设计为50Ω，与射频系统的标准特性阻抗匹配
    \item \textbf{功率容量}：能够承受发射机的最大输出功率，需要良好的散热设计
    \item \textbf{低驻波比}：自身驻波比极低，不影响测试精度
    \item \textbf{频率范围}：覆盖特定的工作频率范围
\end{itemize}

\textbf{阻抗匹配}是假负载最基本的属性。绝大多数通信、业余无线电和测量系统以 $50\,\Omega$ 为标准特性阻抗，有线电视与部分视频系统采用 $75\,\Omega$，两者不可混用：以 $75\,\Omega$ 负载接入 $50\,\Omega$ 系统时，由前述公式可得 $|\Gamma| = 0.2$，$VSWR = 1.5$，虽不至损坏设备，但功率测量将出现约 $4\%$ 的偏差，用于精密测量显然不合适。阻抗的标称值指的是射频阻抗，用万用表测得的直流电阻只能作为粗略核对，对于内部含有隔直电容或分压网络的负载，直流电阻甚至不等于 $50\,\Omega$。

\textbf{功率容量}须区分连续波额定功率与短时（峰值）承受能力。前者由散热能力决定，是热平衡问题；后者由电阻体的瞬时耐压和瞬时能量承受能力决定，是电与介质强度问题。对脉冲或断续工作的场合，可按占空比估算平均功率
\[
P_{av} = P_{pk} \cdot D, \qquad D = \frac{\tau}{T},
\]
式中 $P_{pk}$ 为脉冲峰值功率，$\tau$ 为脉冲宽度，$T$ 为重复周期。需要强调的是，平均功率满足要求并不等于安全，还须核对峰值功率不超过电阻体的瞬时耐受值。此外，许多小型干式负载给出的是"额定功率可持续 X 分钟"的形式，超时使用同样会过热损坏，测试大功率设备时应安排间歇并注意外壳温度。

\textbf{低驻波比}保证假负载本身不引入附加误差。若负载驻波比偏高，部分功率被反射回发射机，功率计读数将随馈线长度变化而起伏，出现"换一根电缆读数就变"的现象；同时反射波与入射波在馈线上形成驻波，可能使某些位置电压过高。因此用于功率测量的负载应选驻波比尽量低的产品，通常测量级负载在其标称频段内可做到较低数值，而以"吸收功率"为主要目的的大功率干式负载指标相对宽松，两者不应混用。

\textbf{频率范围}反映负载保持低驻波比的频段。低频段（中短波）实现容易，单只无感电阻即可胜任；进入甚高频、超高频后，必须采用同轴结构、厚膜片状电阻并严格控制引线长度，才能把寄生参数的影响压到可接受的程度。选用时应确认负载的频率上限覆盖被测设备的工作频率，并留出余量以便测量谐波——测量二次、三次谐波时，负载在谐波频率处的匹配同样影响读数准确性。

\section{应用领域}
\label{sec:dummy_load_applications}
\begin{itemize}
    \item \textbf{设备测试}：测试发射机的输出功率、频率稳定性和调制质量
    \item \textbf{维修调试}：在维修和调试过程中安全地测试设备，避免辐射干扰
    \item \textbf{功率测量}：配合功率计测量发射机的实际输出功率
    \item \textbf{天线替换}：在不希望信号辐射时（如室内测试）替代天线使用
\end{itemize}

\textbf{设备测试}是假负载最主要的用途。发射机的额定输出功率、功率稳定性、载波频率准确度与频率漂移、调制度或频偏、调制线性与失真、谐波与杂散辐射抑制等项目，都要求发射机在标准负载下连续工作，测得的数据才能与技术条件逐条对照。整机定型试验与例行试验中，还需在假负载上做长时间连续发射的老化与温升试验，观察功率下降、频率漂移与保护电路动作情况，这类试验若用实际天线进行既不合法也无法控制条件。

\textbf{维修调试}阶段，假负载是保护被修设备的必备工具。检修末级功放、调整匹配网络、更换功率管后的试机，都应先接假负载，在低功率下确认工作正常后再逐步升到额定功率，最后才转接天线。这样即使匹配网络调整不当或存在自激，损失也仅限于负载发热，不至于连带损坏刚更换的功率器件。同时，维修间内往往同时进行多台设备的检修与测量，用假负载可避免相互干扰造成的误判。

\textbf{功率测量}通常并非把功率计直接接在发射机输出端——普通功率计的输入承受能力远小于发射机输出。规范的连接方式有两种。一是"通过式"接法：发射机经通过式功率计（内含定向耦合器）后接假负载，功率计同时给出正向与反射功率，反射功率读数还可用于核对负载状态。二是"取样"接法：发射机输出接定向耦合器或大功率衰减器的输入端，主路送假负载吸收，耦合端或衰减输出端接功率计或频谱分析仪，被测功率按
\[
P(\mathrm{dBm}) = P_{meas}(\mathrm{dBm}) + A(\mathrm{dB})
\]
换算，其中 $A$ 为耦合度或衰减量，必须使用该器件在测试频率上的实际校准值，否则将直接带来系统误差。在需要频繁切换"天线—假负载"或"多台设备—一台仪表"的场合，可加装同轴开关；使用同轴开关时务必遵守"先停发射、后切换"的原则，带功率热切换会在触点间打火，既烧蚀触点，也可能在切换瞬间造成末级开路。

\textbf{天线替换}用于室内测试、生产检验和教学演示等不希望或不允许辐射的场合。需要注意的是，假负载不能完全消除辐射：负载外壳、接头缝隙与馈线屏蔽层上的漏泄电流仍会产生一定的近场辐射，因此高灵敏度接收测试时应配合屏蔽室或屏蔽箱。另一方面，假负载不能替代天线用于检验天馈系统本身的性能，天线的驻波、方向图与增益必须在实际安装状态下测量。

\section{类型与选择}
\label{sec:dummy_load_types}

假负载根据功率容量和频率范围分为多种类型：

\begin{itemize}
    \item \textbf{低功率假负载}：用于手持对讲机等低功率设备，功率容量通常在几十瓦以下
    \item \textbf{中功率假负载}：用于电台设备，功率容量可达数百瓦
    \item \textbf{高功率假负载}：用于广播发射机等大功率设备，功率容量可达数千瓦甚至更高
    \item \textbf{宽带假负载}：覆盖较宽的频率范围，适用于多波段设备测试
\end{itemize}

按结构与散热方式，假负载还可分为若干形式。\textbf{干式自然冷却}负载靠电阻体、金属外壳与散热片向空气散热，结构简单、无需附件，多见于小功率与中功率产品；\textbf{强迫风冷}负载加装风扇，在同样体积下可显著提高连续功率，但依赖供电，风扇停转或进风口堵塞会迅速过热，故大多设有温度开关或风压联锁。\textbf{油浸式}负载把电阻体浸在绝缘冷却油中，利用油的热容与对流提高短时承载能力，兼有良好的绝缘性能，曾广泛用于短波与广播发射机；使用时须注意油面高度、密封与散热器清洁，切勿在缺油状态下加功率。\textbf{水冷式}负载以循环水带走热量，适用于千瓦以上乃至更大功率的连续消耗，需要配套水泵、热交换器与流量、水温保护。此外，按电气结构可分为集中电阻式（单只或少数几只无感电阻）与同轴分布式（厚膜电阻片安置于渐变同轴结构中），后者频率上限高得多，是微波负载的主流形式。下表概括了各类形式的适用范围。

\begin{table}[htbp]
\centering
\caption{假负载按散热与结构形式的分类概要}
\begin{tabular}{|p{2.4cm}|p{3.0cm}|p{3.2cm}|p{4.0cm}|}
\hline
\textbf{形式} & \textbf{功率量级} & \textbf{散热方式} & \textbf{适用场合与注意事项} \\
\hline
干式自然冷却 & 瓦级至百瓦级 & 外壳、散热片自然对流 & 便携测试、维修间常备；注意连续工作时间限制 \\
\hline
强迫风冷 & 百瓦级至千瓦级 & 风扇强迫对流 & 台架试验与生产测试；须保证风路畅通与风扇可靠 \\
\hline
油浸式 & 百瓦级以上 & 冷却油对流兼绝缘 & 短波与广播发射机试验；注意油面、密封与散热器清洁 \\
\hline
水冷式 & 千瓦级以上 & 循环水与热交换器 & 大功率长时间连续消耗；须设流量与水温保护 \\
\hline
同轴分布式 & 通常为中小功率 & 依结构而定 & 甚高频以上及宽带测量；频率上限高，驻波比低 \\
\hline
\end{tabular}
\end{table}

选择假负载时需注意：
\begin{itemize}
    \item 功率容量应大于被测设备的最大输出功率
    \item 工作频率范围应覆盖被测设备的工作频段
    \item 阻抗必须与被测设备匹配（通常为50Ω）
\end{itemize}

在上述三条基本要求之外，实际选型还应逐项核对以下指标。\textbf{功率余量}：一般取被测设备最大输出功率的 1.5 倍以上，若需长时间连续发射或环境温度较高，余量应更大；同时核对手册给出的额定条件（环境温度、是否带风扇、允许连续时间）是否与实际使用情况相符。\textbf{驻波比与频率的对应关系}：应查看整个使用频段的驻波比曲线或分段指标，而不是只看最好点的数值；用于功率测量时对该指标要求更严。\textbf{接口形式与耐功率}：接头应与被测设备一致，且其本身的功率与频率能力不构成瓶颈。\textbf{散热与安装方式}：需要留出足够的通风空间，负载表面温度较高时应考虑防烫与防止烤伤邻近线缆。\textbf{附加功能}：部分负载带有取样口或内置检波输出，可直接配合功率计或示波器使用，便于测量包络与调制波形，但取样口的耦合度需以校准值为准。最后，负载也是需要维护的器件：长期使用后应定期用驻波表或网络分析仪核对其驻波比，用万用表核对直流电阻，发现变值、变色或有焦糊气味应立即停用。

\section{自制假负载}
\label{sec:diy_dummy_load}

对于低功率应用，假负载可以自制。自制假负载需要注意以下几点：

自制的可行性建立在两个前提上：功率不大，因而散热和耐压容易解决；频率不高，因而寄生参数的影响可以忽略。业余条件下制作短波、中波频段的十瓦级负载并不困难，用于日常检修、发射机试机与调制观察完全够用；但若试图自制甚高频以上或百瓦以上的负载，则元件选择、结构布局和散热都将成为专门问题，其性能通常难以与商品负载相比。

\subsection{电阻选择}

\begin{itemize}
    \item \textbf{电阻值}：必须为50Ω，与射频系统特性阻抗匹配
    \item \textbf{功率容量}：电阻的额定功率应大于被测设备的最大输出功率，建议取2-3倍余量
    \item \textbf{电阻类型}：应使用非电感型电阻，如炭膜电阻、金属膜电阻或水泥电阻；绕线电阻具有电感特性，不适合高频应用
    \item \textbf{频率特性}：高频应用时需考虑电阻的寄生参数，必要时使用多个电阻并联或专用射频电阻
\end{itemize}

电阻类型的选择是自制成败的关键。线绕电阻虽然功率大、价格低，但绕线本身就是一个电感，其感抗随频率线性增长，$X_L = 2\pi f L$，即使只有 $1\ \mu\mathrm{H}$ 的电感量，在 10~MHz 处感抗已达约 $63\,\Omega$，与 $50\,\Omega$ 的电阻分量相当，负载早已不成其为负载。所谓"水泥电阻"多数内部也是线绕结构，只有明确标注无感（非电感型）绕法的才可勉强用于中短波低频端。相对而言，碳膜、金属膜、金属氧化膜电阻的寄生电感小得多，是较合适的选择；专用射频负载电阻（厚膜片式、带法兰便于安装散热）性能最好。

采用多只电阻并联既能提高功率容量，又能降低总的寄生电感，是自制的常用手法。$n$ 只阻值为 $R$ 的电阻并联时
\[
R_p = \frac{R}{n}, \qquad P_{total} = n P_1,
\]
例如四只 $200\,\Omega$ 的电阻并联得到 $50\,\Omega$，总功率为单只的四倍。并联时应把电阻按圆周对称排列在两块导体（如两片圆形铜箔或负载筒的中心导体与外壳）之间，引线尽可能短且长度一致，使电流均匀分配、寄生电感相互抵消一部分。串联同样可行（如两只 $25\,\Omega$ 串联），但串联会累加引线电感，高频特性不如并联。此外还须注意电阻的耐压：单只电阻的最大工作电压限制往往先于功率限制被突破，可用 $U = \sqrt{P R}$ 估算并留出余量。

\subsection{电路结构}

常用的自制假负载电路有两种：

\begin{enumerate}
    \item \textbf{单电阻型}：直接使用一个50Ω电阻，结构简单，适用于频率较低（如中波）或对驻波比要求不高的场合
    
    \item \textbf{π型衰减器}：由三个电阻组成π型网络，可实现更好的阻抗匹配和更宽的频率范围，适合高频应用。典型电路为：两个串联电阻各约16Ω，一个并联电阻约110Ω（根据所需衰减量调整）
\end{enumerate}

单电阻型的电气模型最简单，其等效电路为电阻 $R$ 与分布电容 $C_p$ 并联、再与引线电感 $L_s$ 串联。频率升高时，$C_p$ 的旁路作用使阻抗下降并呈容性，$L_s$ 的作用使阻抗呈感性，两者在某一频率附近还会形成谐振，负载特性急剧恶化。因此单电阻型的可用上限主要取决于结构的紧凑程度，缩短引线、减小对地面积是提高上限的主要办法。

采用电阻性衰减网络的做法，其价值在于用衰减换取匹配：网络输入端的匹配状况主要由网络自身的电阻决定，对输出端所接负载的变化不敏感，衰减量越大这种"隔离"作用越强。同时衰减器还能把过大的功率降到测量仪器可接受的电平，因此测量系统中常把衰减器与负载配合使用。对称电阻衰减器的设计可按下式计算。令所需衰减量为 $A$（dB），定义电压比
\[
K = 10^{A/20},
\]
则 T 型网络（两个串联臂、一个并联臂）的电阻值为
\[
R_1 = R_2 = Z_0\,\frac{K-1}{K+1}, \qquad R_3 = \frac{2 Z_0 K}{K^2 - 1};
\]
$\pi$ 型网络（一个串联臂、两个并联臂）的电阻值为
\[
R_a = Z_0\,\frac{K+1}{K-1}, \qquad R_b = \frac{Z_0 (K^2 - 1)}{2K}.
\]
式中 $Z_0$ 为系统特性阻抗，一般取 $50\,\Omega$。可见网络中各电阻的具体数值完全由所需衰减量决定，上面条目中给出的数值应按实际所选的衰减量重新计算核对。设计时还须注意功率分配：靠近输入端的电阻承担了绝大部分功率，其额定功率必须按输入功率考虑，而不能按平均分配估算。

\subsection{散热设计}

\begin{itemize}
    \item 电阻工作时会发热，必须确保良好的散热
    \item 可将电阻安装在金属散热片上，或使用金属外壳作为散热器
    \item 高功率应用时需注意通风，避免电阻过热损坏
\end{itemize}

散热设计可直接套用前面给出的热阻关系 $T_j = T_a + P R_{th}$ 来估算。假定所用电阻允许的最高体温为 $T_{j\max}$，环境温度为 $T_a$，则允许的功率为
\[
P_{\max} = \frac{T_{j\max} - T_a}{R_{th}} .
\]
这一式子说明了三条实用结论：其一，减小热阻是提高功率能力的唯一途径，故电阻体应紧贴金属外壳或散热片安装，接触面涂导热硅脂，用螺钉压紧而不是仅靠引线悬空固定；其二，环境温度越高允许功率越低，夏季或机柜内部使用时必须降额；其三，加风扇相当于把 $R_{th(s\text{-}a)}$ 降低数倍，效果往往比增大散热片更明显。工程上通常还按经验取 $50\%$ 左右的降额，即用两倍于实际耗散功率的额定电阻，以兼顾长期可靠性。散热片应竖直安装使空气自然对流顺畅，外壳温度可能很高，应设置防烫标识或护网，并远离电缆与易燃物。

\subsection{连接接口}

\begin{itemize}
    \item 接口应与被测设备匹配，常用的有BNC、SMA、PL-259等
    \item 连接线应使用同轴电缆，长度尽量短以减少损耗
    \item 接口处应确保良好的电气连接和屏蔽
\end{itemize}

接口的选择需同时考虑机械形式、功率与频率能力。BNC 采用卡口连接，装拆方便，适用于小功率与较低频段的测量；SMA 为小型螺纹接头，频率上限高但功率能力有限，且拧紧力矩需适当，过力会损坏内导体；PL-259（UHF 型）在短波与业余通信中广泛使用，结构结实、可通过较大功率，但其内部阻抗不连续，不适合甚高频以上的精密测量；N 型与 7/16 型接头则兼顾了功率与频率性能，多用于基站与专业设备。无论何种接头，都应保证中心导体与外导体的连接可靠、焊接饱满无虚焊，外导体屏蔽层与外壳周向连续搭接，否则不仅漏泄增大，还可能在大功率下打火。同轴电缆应尽量短，一方面减小损耗（损耗会使功率计读数低于发射机实际输出），另一方面减小其上驻波造成的影响。

\subsection{自制示例}

一个简单的低功率假负载制作方案：

\begin{itemize}
    \item \textbf{材料}：50Ω/20W水泥电阻1只、BNC插座1个、金属外壳1个
    \item \textbf{制作步骤}：
    \begin{enumerate}
        \item 将水泥电阻固定在金属外壳内壁作为散热
        \item 将电阻一端连接BNC插座的中心引脚，另一端连接外壳（接地）
        \item 确保外壳良好接地，接口处做好屏蔽
    \end{enumerate}
    \item \textbf{适用范围}：适用于10W以下的短波电台测试，频率范围有限（约1-30MHz）
\end{itemize}

制作时的几点补充：电阻应尽量贴近插座安装，引线越短越好，接地端就近连到插座法兰所在的外壳面上，避免形成大面积回路；固定电阻的螺钉应压紧并涂少量导热硅脂；外壳如需开通风孔，孔径应远小于最高工作频率对应的波长，以保持屏蔽效果。

完工后必须先检验再使用，顺序建议为：第一步用万用表测量中心导体与外壳之间的直流电阻，应接近 $50\,\Omega$，并确认无短路与断路；第二步用驻波表或网络分析仪在使用频段内测量驻波比，记录各频点数值，判断可用频率范围；第三步在低功率（如额定值的十分之一）下试机数分钟，用手背靠近感觉温升是否正常，确认无打火、异味；第四步逐步升到实际使用功率并记录允许的连续工作时间。测量结果应做成简单的标签贴在负载上，注明阻抗、额定功率、允许连续时间与经检验的频率范围，避免日后误用。

\subsection{注意事项}

\begin{itemize}
    \item 自制假负载的驻波比通常不如商业产品，适用于要求不高的测试场合
    \item 高频应用时需注意分布参数的影响，必要时进行阻抗匹配调整
    \item 切勿在超过电阻额定功率的情况下使用，以免造成设备损坏
    \item 自制假负载不建议用于精密测量或大功率设备测试
\end{itemize}

除以上四条，操作安全同样需要强调。假负载在工作时表面温度可能相当高，发射结束后仍会持续散热，切勿立即用手触摸或收进箱包；百瓦级以上设备的负载端射频电压可达上百伏，发射状态下不得触碰接头、拆装电缆或切换同轴开关。使用中应养成"先接负载、后加激励；先停发射、后拆负载"的习惯，任何时候都不要让发射机处于空载或短路状态。此外，长期使用后的自制负载应定期复测直流电阻与驻波比：电阻在反复热循环后可能阻值漂移或引线疲劳断裂，一旦阻值明显偏离 $50\,\Omega$ 就应更换。最后需要明确，自制负载的定位是"安全试机与粗略测量"，凡涉及指标判定、报告出具或设备验收的测量，都应使用经过校准的商品负载与配套仪表。
